% !TeX document-id = {40e7f911-37dd-4856-a048-b5b803f1756a}
% !TeX TXS-program:compile = txs:///pdflatex/[-shell-escape]

\documentclass[11pt]{beamer}
\usepackage[utf8]{inputenc}
\usepackage[T1,T2A]{fontenc}
\usepackage[russian]{babel}
\usepackage{color}
\usepackage{listings}
\usepackage{calc}
\usepackage{graphicx}
\usepackage{epstopdf}
\usepackage{hyperref}
\hypersetup{unicode,colorlinks}
\usepackage{csquotes}
\usepackage{upquote}
\usepackage{cprotect}
\usetheme[progressbar=head,numbering=fraction,block=fill]{metropolis}
\usepackage{dejavu}
\usepackage{etoolbox}
\usepackage{bm}
\usepackage{minted}
%\usepackage[defaultsans]{droidsans}
%\usepackage[defaultmono]{droidmono}
%\makeatletter
%\input droid/t1fdm.fd
%\makeatother

\makeatletter
\DeclareRobustCommand*{\bfseries}{%
    \not@math@alphabet\bfseries\mathbf
    \fontseries\bfdefault\selectfont
    \boldmath
}
\makeatother

\makeatletter
\chardef\straightquote@code=\catcode`'
\chardef\backquote@code=\catcode``
\catcode`'=\active \catcode``=\active
\patchcmd{\@noligs}
{\textasciigrave}
{\fixedtextasciigrave}
{}{}
\newcommand{\fixedtextasciigrave}{%
    \makebox[.5em]{\fontencoding{TS1}\fontfamily{fvs}\selectfont\textasciigrave}% Vera Sans
}
\catcode\lq\'=\straightquote@code
\catcode\lq\`=\backquote@code
\makeatletter

\lstset{
    language=Haskell,
    inputencoding=utf8,
    extendedchars=true,
    breaklines=true,
    escapeinside=!!,
    tabsize=4,
    breakatwhitespace=true,
    keepspaces=true,
    basicstyle=\ttfamily,
    commentstyle=\color{gray}\itshape,
    keywordstyle=
    % keywordstyle=\color{blue}, % TODO https://tex.stackexchange.com/questions/415777/avoid-highlighting-keywords-following-certain-words-in-listings
    % addkeywords={Type}
    % deletekeywords={Monad, String, return, array, show},
    % otherkeywords={::, ->}
}
\lstset{
    literate={а}{{\selectfont\char224}}1
    {б}{{\selectfont\char225}}1
    {в}{{\selectfont\char226}}1
    {г}{{\selectfont\char227}}1
    {д}{{\selectfont\char228}}1
    {е}{{\selectfont\char229}}1
    {ё}{{\"e}}1
    {ж}{{\selectfont\char230}}1
    {з}{{\selectfont\char231}}1
    {и}{{\selectfont\char232}}1
    {й}{{\selectfont\char233}}1
    {к}{{\selectfont\char234}}1
    {л}{{\selectfont\char235}}1
    {м}{{\selectfont\char236}}1
    {н}{{\selectfont\char237}}1
    {о}{{\selectfont\char238}}1
    {п}{{\selectfont\char239}}1
    {р}{{\selectfont\char240}}1
    {с}{{\selectfont\char241}}1
    {т}{{\selectfont\char242}}1
    {у}{{\selectfont\char243}}1
    {ф}{{\selectfont\char244}}1
    {х}{{\selectfont\char245}}1
    {ц}{{\selectfont\char246}}1
    {ч}{{\selectfont\char247}}1
    {ш}{{\selectfont\char248}}1
    {щ}{{\selectfont\char249}}1
    {ъ}{{\selectfont\char250}}1
    {ы}{{\selectfont\char251}}1
    {ь}{{\selectfont\char252}}1
    {э}{{\selectfont\char253}}1
    {ю}{{\selectfont\char254}}1
    {я}{{\selectfont\char255}}1
    {А}{{\selectfont\char192}}1
    {Б}{{\selectfont\char193}}1
    {В}{{\selectfont\char194}}1
    {Г}{{\selectfont\char195}}1
    {Д}{{\selectfont\char196}}1
    {Е}{{\selectfont\char197}}1
    {Ё}{{\"E}}1
    {Ж}{{\selectfont\char198}}1
    {З}{{\selectfont\char199}}1
    {И}{{\selectfont\char200}}1
    {Й}{{\selectfont\char201}}1
    {К}{{\selectfont\char202}}1
    {Л}{{\selectfont\char203}}1
    {М}{{\selectfont\char204}}1
    {Н}{{\selectfont\char205}}1
    {О}{{\selectfont\char206}}1
    {П}{{\selectfont\char207}}1
    {Р}{{\selectfont\char208}}1
    {С}{{\selectfont\char209}}1
    {Т}{{\selectfont\char210}}1
    {У}{{\selectfont\char211}}1
    {Ф}{{\selectfont\char212}}1
    {Х}{{\selectfont\char213}}1
    {Ц}{{\selectfont\char214}}1
    {Ч}{{\selectfont\char215}}1
    {Ш}{{\selectfont\char216}}1
    {Щ}{{\selectfont\char217}}1
    {Ъ}{{\selectfont\char218}}1
    {Ы}{{\selectfont\char219}}1
    {Ь}{{\selectfont\char220}}1
    {Э}{{\selectfont\char221}}1
    {Ю}{{\selectfont\char222}}1
    {Я}{{\selectfont\char223}}1
    {`}{{{\fixedtextasciigrave}}}1
}

% \usemintedstyle{solarized-light}
\newminted[haskell]{haskell}{
    escapeinside=``,
    beameroverlays=true,
    texcomments=true,
    autogobble=true,
    fontsize=\small
}
\newmintinline[haskinline]{haskell}{
    escapeinside=``,
    beameroverlays=true,
    texcomments=true
}

\author{Алексей Романов}
\subtitle{Функциональное программирование на Haskell}
%\logo{}
\institute{МИЭТ}
\subject{Функциональное программирование на Haskell}
%\setbeamercovered{transparent}
%\setbeamertemplate{navigation symbols}{}


\title{Лекция 11: монада IO и взаимодействие с внешним миром}

\begin{document}

\begin{frame}[plain]
  \maketitle
\end{frame}

\begin{frame}[fragile]
  \frametitle{Монада IO}
  \begin{itemize}
    \item Среди многих монад стандартной библиотеки, \haskinline|IO| играет особую роль.
    \item Значение типа \haskinline|IO a| описывает вычисление с результатом типа \haskinline|a| и побочными эффектами.
    \item \haskinline|IO| описывает взаимодействие с внешним миром и работу с изменяемым состоянием.
  \end{itemize}
\end{frame}

\begin{frame}[fragile]
  \frametitle{Программы командной строки}
  \begin{itemize}
    \item Исполняемая программа в Haskell задаётся значением
          \haskinline|main :: IO a| (почти всегда \haskinline|IO ()|).
    \item Это аналог \mintinline{c}|int main(int argc, char** argv)| в C, \mintinline{java}|public static void main(String[] args)| в Java и т.д.
    \item Соответственно, есть доступ к аргументам командной строки и возможность указать код выхода:
    \item \href{https://hackage.haskell.org/package/base/docs/System-Environment.html}{\haskinline{System.Environment.getArgs :: IO [String]}}. Там же есть функции для чтения переменных окружения.
    \item \href{https://hackage.haskell.org/package/base/docs/System-Exit.html}{\haskinline{System.Exit.exitWith :: ExitCode -> IO a}}.
  \end{itemize}
\end{frame}

\begin{frame}[fragile]
  \frametitle{Интерактивные программы}
  \begin{itemize}
    \item Из консоли можно читать текст, вводимый пользователем:
          \begin{haskell}
            getLine ::!\pause! IO String
          \end{haskell}
    \item И писать в неё:
          \begin{haskell}
            putStr, putStrLn ::!\pause! String -> IO ()
            print :: Show a => a -> IO () 
          \end{haskell}
          \pause
    \item Простой пример: читаем любую строку, выводим её длину и повторяем (а на \haskinline|"quit"| выходим):
          \pause
          \begin{haskellsmall}
            main = do
              line <- getLine
              if line /= "quit"
                then do
                  print (length line)
                  main
                else
                  pure ()  -- или exitSuccess
          \end{haskellsmall}
  \end{itemize}
\end{frame}

\begin{frame}[fragile]
  \frametitle{Побег из \haskinline|IO|}
  \begin{itemize}
    \item Если у нас есть значение типа \haskinline|IO a|, можно ли превратить его просто в \haskinline|a|?
          \pause
    \item В обычном коде нет.
          \pause
    \item Внутри \haskinline|do| можно использовать результат действия, но само вычисление остаётся в \haskinline|IO|.
    \item Но есть функция \\
          \href{https://hackage.haskell.org/package/base/docs/System-IO-Unsafe.html#v:unsafePerformIO}{\haskinline{unsafePerformIO :: IO a -> a}}.
    \item Само название говорит о её небезопасности.
    \item Конкретное использование может быть и безопасным, но для этого надо знать довольно много о внутренностях Haskell.
          \pause
    \item Поэтому в рамках этого курса мы её использовать не будем, как и \href{https://hackage.haskell.org/package/base/docs/System-IO-Unsafe.html}{другие \haskinline{unsafe*IO}}.
    \item А есть ещё \href{https://hackage.haskell.org/package/bytestring/docs/Data-ByteString-Internal.html#v:accursedUnutterablePerformIO}{\haskinline{accursedUnutterablePerformIO}}\ldots
  \end{itemize}
\end{frame}

\begin{frame}[fragile]
  \frametitle{Обход с \haskinline|IO|: \haskinline|for| и \haskinline|for_|}
  \begin{itemize}
    \item Как \haskinline|traverse| и \haskinline|traverse_|, но функция передаётся последней:
          \begin{haskellsmall}
            import Data.Traversable (for)
            import Data.Foldable (for_)
          \end{haskellsmall}
          \haskinline|for = flip traverse|, \haskinline|for_ = flip traverse_|.
    \item \haskinline|for| собирает результаты; \haskinline|for_| отбрасывает их, сохраняя эффекты. \haskinline|forM(_)| и \haskinline|mapM(_)| "--- старые имена того же для монад.
          \pause
    \item Так удобно записывать действие из нескольких строк:
          \begin{haskell}
            for_ [1..3] $ \i -> do
              print i
              print (i * i)
          \end{haskell}
          Для каждого элемента печатаем число и его квадрат.
  \end{itemize}
\end{frame}

\begin{frame}[fragile]
  \frametitle{Условное выполнение: \haskinline|when|}
  \begin{itemize}
    \item Функция из упражнения в лекции 9:
          \begin{haskellsmall}
            import Control.Monad (when)

            when :: Applicative f => Bool -> f () -> f ()
          \end{haskellsmall}
    \item При \haskinline|True| выполняет действие, при \haskinline|False| возвращает \haskinline|pure ()| без его эффектов.
          \pause
    \item Например, условный диагностический вывод:
          \begin{haskellsmall}
            report :: Bool -> IO ()
            report debug = do
              when debug (putStrLn "Debug mode")
              putStrLn "Running"
          \end{haskellsmall}
          \haskinline|report False| печатает только \haskinline|"Running"|.
  \end{itemize}
\end{frame}

\begin{frame}[fragile]
  \frametitle{Ловушки консоли: буферизация}
  \begin{itemize}
    \item В скомпилированной программе \haskinline|stdout| буферизуется построчно (а при перенаправлении в файл или канал "--- блоками). Поэтому
          \begin{haskell}
            putStr "Введите имя: "
            name <- getLine
          \end{haskell}
          не покажет приглашение, пока пользователь не нажмёт Enter.
          \pause
    \item Решения: \haskinline|hFlush stdout| после \haskinline|putStr| или \haskinline|hSetBuffering stdout NoBuffering| в начале \haskinline|main| (модуль \href{https://hackage.haskell.org/package/base/docs/System-IO.html}{\haskinline{System.IO}}).
    \item В GHCi и \ghcinline|stack repl| этого не видно: там буферизация отключена. Проверяйте \ghcinline|stack run|.
  \end{itemize}
\end{frame}

\begin{frame}[fragile]
  \frametitle{Ловушки консоли: кодировки}
  \begin{itemize}
    \item Кодировка \haskinline|stdin|/\haskinline|stdout| и \haskinline|readFile|/\haskinline|writeFile| определяется окружением; UTF-8 не гарантирована.
    \item Под Windows при неподходящей кодировке программа может упасть или вывести крякозябры.
          \pause
    \item Решение: \ghcinline|chcp 65001| перед запуском и/или в программе
          \begin{haskell}
            hSetEncoding stdout utf8
            hSetEncoding stdin utf8
          \end{haskell}
    \item Для файлов "--- \haskinline|hSetEncoding h utf8| после \haskinline|openFile| (или \haskinline|withFile|).
  \end{itemize}
\end{frame}

\begin{frame}[fragile]
  \frametitle{Изменяемые переменные: \haskinline|IORef|}
  \begin{itemize}
    \item Модуль \href{https://hackage.haskell.org/package/base/docs/Data-IORef.html}{\haskinline{Data.IORef}}:
          \begin{haskell}
            newIORef :: a -> IO (IORef a)
            readIORef :: IORef a -> IO a
            writeIORef :: IORef a -> a -> IO ()
            modifyIORef' :: IORef a -> (a -> a) -> IO ()
          \end{haskell}
    \item Используйте \haskinline|modifyIORef'| (строгий), а не \haskinline|modifyIORef|: последний накапливает отложенные вычисления ровно как \haskinline|foldl| в лекции 8.
    \item Для нескольких потоков нужны \haskinline|atomicModifyIORef'|, \haskinline|MVar| или \haskinline|STM| (лекция 13).
  \end{itemize}
\end{frame}

\begin{frame}[fragile]
  \frametitle{Исключения (кратко)}
  \begin{itemize}
    \item Модуль \href{https://hackage.haskell.org/package/base/docs/Control-Exception.html}{\haskinline{Control.Exception}}:
          \begin{haskellsmall}
            throwIO :: Exception e => e -> IO a
            catch   :: Exception e => IO a -> (e -> IO a) -> IO a
            try     :: Exception e => IO a -> IO (Either e a)
            evaluate :: a -> IO a
            bracket :: IO a -> (a -> IO b) -> (a -> IO c) -> IO c
          \end{haskellsmall}
    \item Нужно определить, какие исключения ловить:
          \begin{haskellsmall}
            import System.IO (readFile')

            try @IOException (readFile' path)
          \end{haskellsmall}
          Ловить всё (\haskinline|SomeException|) "--- обычно плохая идея.
    \item \haskinline|error| и неполные сопоставления тоже исключения (\haskinline|ErrorCall|, \haskinline|PatternMatchFail|), но лучше их не ловить, а не допускать.
  \end{itemize}
\end{frame}

\begin{frame}[fragile]
  \frametitle{Исключения и ленивость}
  \begin{itemize}
    \item \haskinline|pure| не вычисляет результат:
          \begin{haskellsmall}
            try @ArithException (pure (div (1 :: Int) 0))
          \end{haskellsmall}
          Получим \haskinline|Right| с отложенным вычислением; ошибка возникнет уже вне \haskinline|try|.
          \pause
    \item Чтобы вычислить результат внутри \haskinline|try|:
          \begin{haskellsmall}
            try @ArithException (evaluate (div (1 :: Int) 0))
          \end{haskellsmall}
          Здесь получим \haskinline|Left DivideByZero|.
    \item \haskinline|evaluate| вычисляет только до \emph{WHNF} (лекция 8).
    \item Ошибки ленивого \haskinline|readFile| тоже могут возникнуть вне \haskinline|try|; \haskinline|readFile'| читает файл сразу.
  \end{itemize}
\end{frame}

\begin{frame}[fragile]
  \frametitle{Работа с файлами}
  \begin{itemize}
    \item Модуль \href{https://hackage.haskell.org/package/base/docs/System-IO.html}{\haskinline|System.IO|}.
    \item Напомню, \haskinline|type FilePath = String|.
    \item Простейшие операции это чтение и запись в файл:
          \begin{haskell}
            readFile :: FilePath -> IO String
            writeFile, appendFile :: 
              FilePath -> String -> IO ()
          \end{haskell}
    \item
          Или можно открыть файл один раз, получить \haskinline|Handle| (дескриптор файла) и работать с ним:
          \begin{haskell}
            openFile, openBinaryFile :: 
              FilePath -> IOMode -> IO Handle 
            hPutStrLn :: Handle -> String -> IO ()
            hGetLine, hGetContents :: Handle -> IO String
          \end{haskell}
          \pause
    \item
          Функции работы с консолью сводятся к этим:
          \begin{haskell}
            getLine = hGetLine stdin
          \end{haskell}
  \end{itemize}
\end{frame}

\begin{frame}[fragile]
  \frametitle{Работа с файлами (2)}
  \begin{itemize}
    \item В чём проблема с кодом вроде
          \begin{haskell}
            do
              file <- openFile path ReadWriteMode
              ...
              hClose file
          \end{haskell}
          \pause
    \item Что, если на каком-то из шагов случится исключение? Мы остаёмся с открытым файлом.
    \item Если это случится много раз, программа вылетит.
          \pause
    \item Чтобы закрыть файл и при исключении:
          \begin{haskell}
            withFile path ReadWriteMode $ \file -> do
              ...
          \end{haskell}
    \item Это частный случай функции \haskinline|bracket|.
  \end{itemize}
\end{frame}

\begin{frame}[fragile]
  \frametitle{Ленивый ввод-вывод}
  \begin{itemize}
    \item \haskinline|hGetContents| возвращает строку: файл читается по мере её вычисления.
    \item С одной стороны, это хорошо: не нужно явно указывать, сколько читать.
    \item Но попытка дочитать строку после закрытия файла вызывает ошибку:
          \begin{haskellsmall}
            wrong = do
              fileData <-
                withFile "test.txt" ReadMode hGetContents
              putStr fileData
          \end{haskellsmall}
    \item Нужно использовать данные внутри \haskinline|withFile|:
          \begin{haskellsmall}
            right = withFile "test.txt" ReadMode $ \file -> do
              fileData <- hGetContents file
              putStr fileData
          \end{haskellsmall}
  \end{itemize}
\end{frame}

\begin{frame}[fragile]
  \frametitle{Строгий ввод-вывод}
  \begin{itemize}
    \item С GHC 9.0 в \haskinline|System.IO| есть строгие варианты \haskinline|readFile'|, \haskinline|hGetContents'|, \haskinline|getContents'|: они читают всё сразу, и проблем с закрытием файла у них нет.
    \item Обычно лучше использовать их или строгие варианты \haskinline|Text|/\haskinline|ByteString| со следующего слайда. Ленивые варианты \haskinline|readFile| читают по требованию.
    \item Ленивый ввод-вывод остаётся полезным для потоковой обработки больших файлов, но для этого есть и более предсказуемые библиотеки (\hackage{conduit}, \hackage{streamly}).
  \end{itemize}
\end{frame}

\begin{frame}[fragile]
  \frametitle{\haskinline|ByteString| и \haskinline|Text|}
  \begin{itemize}
    \item Мы знаем, что \haskinline|String| занимает очень много места в памяти, что ухудшает и время работы.
    \item Вместо него есть \href{https://hackage.haskell.org/package/bytestring/docs/Data-ByteString.html}{\haskinline{Data.ByteString.ByteString}}, по сути представляющий собой массив байтов, и \href{https://hackage.haskell.org/package/bytestring/docs/Data-ByteString-Lazy.html}{\haskinline{Data.ByteString.Lazy.ByteString}} как ленивый список таких массивов.
    \item Многие функции для работы с ними называются так же, как для \haskinline|String|, но живут в соответствующих модулях (поэтому используется \haskinline|import qualified|).
    \item Они годятся для бинарных данных (или ASCII, с помощью \href{https://hackage.haskell.org/package/bytestring/docs/Data-ByteString-Char8.html}{\haskinline{Data.ByteString}}\haskinline|[|\href{https://hackage.haskell.org/package/bytestring/docs/Data-ByteString-Lazy-Char8.html}{\haskinline{.Lazy}}\haskinline|].Char8|).
    \item Для текста аналогичный пакет \haskinline|text| и типы в модулях \href{https://hackage.haskell.org/package/text/docs/Data-Text.html}{\haskinline{Data.Text}}\haskinline|[|\href{https://hackage.haskell.org/package/text/docs/Data-Text-Lazy.html}{\haskinline{.Lazy}}\haskinline|]|.
  \end{itemize}
\end{frame}

\begin{frame}[fragile]
  \frametitle{Текст и байты: UTF-8}
  \begin{itemize}
    \item \haskinline|ByteString| содержит байты, \haskinline|Text| "--- текст. Для перехода между ними нужно выбрать кодировку.
    \item \haskinline|Data.ByteString.Char8| \emph{не декодирует UTF-8}: каждый байт интерпретируется как отдельный символ.
          \pause
    \item Для UTF-8 есть модуль \href{https://hackage.haskell.org/package/text/docs/Data-Text-Encoding.html}{\haskinline{Data.Text.Encoding}}:
          \begin{haskellsmall}
            decodeUtf8' :: ByteString ->
              Either UnicodeException Text
            encodeUtf8 :: Text -> ByteString
          \end{haskellsmall}
    \item \haskinline|decodeUtf8'| возвращает \haskinline|Left| при некорректной последовательности байтов; \haskinline|Right| содержит декодированный текст.
  \end{itemize}
\end{frame}

\begin{frame}[fragile]
  \frametitle{Случайные значения}
  \begin{itemize}
    \item Модуль \href{https://hackage.haskell.org/package/random/docs/System-Random.html}{\haskinline{System.Random}} (пакет \hackage{random}). Класс \haskinline|RandomGen| описывает чистые генераторы псевдослучайных чисел. Например:
          \begin{haskellsmall}
            genWord64 :: RandomGen g => g -> (Word64, g)
          \end{haskellsmall}
    \item Получаем число и \emph{новое состояние} генератора. \haskinline|Word64| "--- 64-битное беззнаковое целое из \href{https://hackage.haskell.org/package/base/docs/Data-Word.html}{\haskinline{Data.Word}}.
    \item Стандартный генератор \haskinline|StdGen| имеет экземпляр \haskinline|RandomGen|. Его можно создать из начального значения (\emph{seed}):
          \begin{haskellsmall}
            mkStdGen :: Int -> StdGen
          \end{haskellsmall}
  \end{itemize}
\end{frame}

\begin{frame}[fragile]
  \frametitle{Случайные значения: тип результата}
  \begin{itemize}
    \item \haskinline|Uniform a|: равномерное распределение по всем значениям типа \haskinline|a| (например, \haskinline|Bool| или \haskinline|Int|).
    \item \haskinline|UniformRange a|: равномерное распределение в заданном диапазоне (например, для \haskinline|Int|, \haskinline|Integer| или \haskinline|Double|).
    \item В \haskinline|System.Random| есть чистые функции:
          \begin{haskellsmall}
            uniform :: (Uniform a, RandomGen g) =>
              g -> (a, g)
            uniformR :: (UniformRange a, RandomGen g) =>
              (a, a) -> g -> (a, g)
          \end{haskellsmall}
    \item Результат и новое состояние генератора "--- знакомая схема \haskinline|State|.
  \end{itemize}
\end{frame}

\begin{frame}[fragile]
  \frametitle{Случайные значения: воспроизводимость}
  \begin{itemize}
    \item Передаём новое состояние в следующий вызов:
          \begin{haskellsmall}
            g0 = mkStdGen 42
            (n1 :: Int, g1) = uniformR (1, 100) g0
            (n2 :: Int, g2) = uniformR (1, 100) g1
          \end{haskellsmall}
          \begin{haskellsmall}
            ghci> (n1, n2)
            (49,24)
          \end{haskellsmall}
    \item Если снова передать \haskinline|g0|, \pause получим \haskinline|(49,49)|.
    \item При той же версии \haskinline|random| повторный запуск с тем же seed даёт ту же последовательность.
  \end{itemize}
\end{frame}

\begin{frame}[fragile]
  \frametitle{Случайные значения в IO}
  \begin{itemize}
    \item В \href{https://hackage.haskell.org/package/random/docs/System-Random-Stateful.html}{\haskinline{System.Random.Stateful}} генераторы сами хранят изменяемое состояние:
          \begin{haskellsmall}
            uniformM :: (Uniform a, StatefulGen g m) =>
              g -> m a
            uniformRM :: (UniformRange a, StatefulGen g m) =>
              (a, a) -> g -> m a
          \end{haskellsmall}
          \haskinline|StatefulGen g m| связывает генератор \haskinline|g| с монадой \haskinline|m|.
    \item Глобальный генератор обновляется атомарно:
          \begin{haskellsmall}
            globalStdGen :: AtomicGenM StdGen
          \end{haskellsmall}
    \item Например, бросок кубика:
          \begin{haskellsmall}
            rollDie :: IO Int
            rollDie = uniformRM (1, 6) globalStdGen
          \end{haskellsmall}
          Каждый вызов обновляет состояние генератора.
  \end{itemize}
\end{frame}

\begin{frame}[fragile]
  \frametitle{Из чистого генератора в IO}
  \begin{itemize}
    \item \haskinline|AtomicGenM| "--- обёртка над \haskinline|IORef|:
          \begin{haskellsmall}
            newtype AtomicGenM g = AtomicGenM (IORef g)
            atomicModifyIORef' :: IORef g -> (g -> (g, a)) -> IO a
          \end{haskellsmall}
    \item Реализуем библиотечную функцию (тип для \haskinline|IO|):
          \begin{overprint}[\linewidth]
            \onslide<1>
              \begin{haskellsmall}
                applyAtomicGen :: (g -> (a, g)) -> AtomicGenM g -> IO a
                applyAtomicGen step (AtomicGenM ref) = ???
              \end{haskellsmall}
            \onslide<2->
              \begin{haskellsmall}
                applyAtomicGen :: (g -> (a, g)) -> AtomicGenM g -> IO a
                applyAtomicGen step (AtomicGenM ref) =
                  atomicModifyIORef' ref $ \g ->
                    let (x, g') = step g
                    in (g', x)
              \end{haskellsmall}
          \end{overprint}
          \pause
    \item Тот же бросок кубика через \haskinline|applyAtomicGen|:
          \begin{haskellsmall}
            rollDie :: IO Int
            rollDie = applyAtomicGen (uniformR (1, 6)) globalStdGen
          \end{haskellsmall}
  \end{itemize}
\end{frame}

\begin{frame}[fragile]
  \frametitle{Дополнительное чтение}
  \begin{itemize}
    \item \href{https://www.engineering.upenn.edu/~cis1940/spring13/lectures/08-IO.html}{CIS 194: IO}: действия как описания вычислений и связь с чистотой.
    \item \href{https://wiki.haskell.org/IO_inside}{IO inside} (Haskell Wiki): как устроена монада \haskinline|IO|.
    \item \href{https://www.snoyman.com/blog/2016/12/beware-of-readfile/}{Beware of readFile}: ловушки ленивого ввода-вывода и кодировок.
    \item \href{https://academy.fpblock.com/haskell/tutorial/exceptions/}{Safe exception handling}: исключения и освобождение ресурсов; примеры с \haskinline|UnliftIO.Exception|/\haskinline|RIO|.
    \item Simon Peyton Jones: \href{https://www.microsoft.com/en-us/research/publication/tackling-awkward-squad-monadic-inputoutput-concurrency-exceptions-foreign-language-calls-haskell/}{Tackling the Awkward Squad} (\emph{углублённо}): \haskinline|IO|, исключения, потоки и FFI; часть API устарела (2001).
  \end{itemize}
\end{frame}

\end{document}
